\slajdid{09_2-s052}
\begin{frame}[label=09_2-s052]
Lečap može da bude destruktivan za komponentu, ali u mnogim slučajevima disipacija nije dovoljna za destrukciju zbog ograničenog strujnog kapaciteta izvora za napajanje. Ali će baš iz tog razloga napon napajanja pasti i onemogućiti ispravno funkcionisanje komponente.\par\bigskip
\alert{Ako se lečap pojavio, da bi nestao mora da se ukloni napon napajanja.}\par\bigskip
Pojava lečapa najčešće nastaje kod delova kola koja su povezana na „spojni svet“, kada prilikom nestajanja ili uspostavljanja napona napon na izlaznim priključcima naglo raste. Zbog toga se često kod ovih komponenti povećava vreme uspona i pada signala, odnosno smanjuje brzina porasta signala na izlazu. Druga tehnika je da se oblasti P i N tranzistora okruže prstenovima koji će biti direktno povezani na napajanje odnosno na masu kako bi se efektivno smanjio uticaj jednog tranzistora na drugi.
\end{frame}
